\documentclass[11pt]{article}

\usepackage[margin=1in]{geometry}
\usepackage{amsmath,amssymb,amsthm}
\usepackage{bm}
\usepackage{algorithm}
\usepackage{algpseudocode}
\usepackage{hyperref}
\usepackage{cleveref}
\usepackage{booktabs}

\newcommand{\R}{\mathbb{R}}
\newcommand{\tr}{\operatorname{tr}}
\newcommand{\diag}{\operatorname{diag}}
\newcommand{\ve}{\operatorname{vec}}
\newcommand{\dd}{\mathrm{d}}
\newcommand{\bx}{\bm{x}}
\newcommand{\by}{\bm{y}}
\newcommand{\bs}{\bm{s}}
\newcommand{\bc}{\bm{c}}
\newcommand{\bmu}{\bm{\mu}}
\newcommand{\bep}{\bm{\varepsilon}}
\newcommand{\bla}{\bm{\lambda}}
\newcommand{\sstar}{\bs^{\star}}
\newcommand{\Snew}{\bm{\Sigma}_{\mathrm{new}}}
\newcommand{\FIM}{\bm{I}_{\mathrm{F}}}

\theoremstyle{definition}
\newtheorem{definition}{Definition}
\newtheorem{proposition}{Proposition}
\newtheorem{remark}{Remark}

\title{Posterior Covariance Minimization for Adaptive Sensor Selection\\[4pt]
       \large Bilevel Formulation and End-to-End Backpropagation}
\author{}
\date{}

\begin{document}
\maketitle

\begin{abstract}
We present a bilevel optimization framework for sensor design in which the
inner problem selects sensor parameters to minimize the expected posterior
covariance of a Bayesian state estimator, and the outer problem optimizes
the physical sensor design to minimize long-horizon estimation error.  We
derive the closed-form sensor selection policy under a linearized Gaussian
measurement model, show that the expected posterior covariance is
independent of the realized measurement, and develop an efficient
backpropagation scheme through the inner argmin via the Implicit Function
Theorem.  We give explicit expressions for all required derivatives,
including a novel analytical directional-derivative formulation that
eliminates nested automatic differentiation.
\end{abstract}

%======================================================================
\section{Problem Setup}
\label{sec:setup}
%======================================================================

Consider a hidden state $\bx \in \R^{d_x}$ observed through a nonlinear
sensor model
\begin{equation}
  \by = f(\bx, \bs, \bc) + \bep, \qquad
  \bep \sim \mathcal{N}(\bm{0},\, \sigma^2 \bm{I}_{d_y}),
  \label{eq:obs-model}
\end{equation}
where $\bs \in \R^{d_s}$ are \emph{sensor parameters} (e.g., input
phases) chosen adaptively at each time step, and $\bc \in \R^{d_c}$ are
\emph{sensor design parameters} (e.g., a photonic device geometry)
optimized offline.

A Bayesian estimator maintains a belief $(\bmu, \bm{\Sigma})$ over $\bx$.
At each step, we (i) select $\bs$ to maximize estimation quality, (ii)
acquire a measurement $\by$, and (iii) update the belief via the Extended
Kalman Filter (EKF).

%======================================================================
\section{Expected Posterior Covariance}
\label{sec:posterior-cov}
%======================================================================

\subsection{EKF Update}

Linearizing \eqref{eq:obs-model} around the current mean $\bmu$ gives
\begin{equation}
  \bm{F} = \frac{\partial f}{\partial \bx}\bigg|_{\bmu, \bs, \bc}
  \in \R^{d_y \times d_x}.
\end{equation}
The standard EKF measurement update is:
\begin{align}
  \bm{S}     &= \bm{F}\bm{\Sigma}\bm{F}^\top + \sigma^2 \bm{I}_{d_y},
  \label{eq:innov-cov} \\
  \bm{K}     &= \bm{\Sigma}\bm{F}^\top \bm{S}^{-1},
  \label{eq:kalman-gain} \\
  \bmu_{\mathrm{new}} &= \bmu + \bm{K}(\by - \hat{\by}),
  \label{eq:mean-update} \\
  \Snew      &= (\bm{I}_{d_x} - \bm{K}\bm{F})\,\bm{\Sigma}.
  \label{eq:cov-update}
\end{align}

\subsection{Key Observation: Independence from $\by$}

\begin{proposition}
\label{prop:y-indep}
In the linearized Gaussian model, the posterior covariance $\Snew$ depends
on $(\bmu, \bm{\Sigma}, \bs, \bc)$ but \textbf{not} on the realized
measurement~$\by$.
\end{proposition}

\begin{proof}
The Kalman gain $\bm{K}$ in \eqref{eq:kalman-gain} depends only on
$\bm{F}$, $\bm{\Sigma}$, and $\sigma^2$, all of which are determined
before $\by$ is observed.  Since $\Snew = (\bm{I} - \bm{K}\bm{F})\bm{\Sigma}$
involves only $\bm{K}$ and $\bm{F}$, it is independent of $\by$.
\end{proof}

Consequently, the \emph{expected} posterior covariance $\mathbb{E}_{\by}[\Snew]$
equals $\Snew$ itself---no integration over $\by$ is required.

\subsection{Woodbury Form}

Applying the matrix inversion lemma to \eqref{eq:cov-update} yields the
compact expression
\begin{equation}
  \boxed{%
    \Snew(\bs) = \bigl(\bm{\Sigma}^{-1}
      + \bm{F}(\bs)^\top \bm{F}(\bs) / \sigma^2\bigr)^{-1}
    = \bigl(\bm{\Sigma}^{-1} + \FIM(\bs)\bigr)^{-1},
  }
  \label{eq:woodbury}
\end{equation}
where $\FIM = \bm{F}^\top\bm{F}/\sigma^2$ is the (Bayesian) Fisher
Information Matrix for a single measurement.  The posterior precision is
the sum of the prior precision and the Fisher information.

%======================================================================
\section{Sensor Selection Policy}
\label{sec:sensor-selection}
%======================================================================

\subsection{A-Optimal Criterion}

We select sensor parameters to minimize the trace of the posterior
covariance (the A-optimality criterion in experimental design):
\begin{equation}
  \sstar(\bmu, \bm{\Sigma}, \bc)
  = \arg\min_{\bs}\;\underbrace{%
    \tr\!\bigl(\Snew(\bs)\bigr) + \alpha_r \|\bs\|^2
  }_{\displaystyle \varphi(\bs;\, \bmu, \bm{\Sigma}, \bc)},
  \label{eq:inner-opt}
\end{equation}
where $\alpha_r > 0$ is a regularization coefficient that keeps the
solution bounded.

\begin{remark}[Comparison with the BFIM criterion]
The commonly used BFIM criterion maximizes $\tr(\FIM) = \|\bm{F}\|_F^2 / \sigma^2$,
which depends only on $\bm{F}$ and $\sigma^2$.  The A-optimal criterion
\eqref{eq:inner-opt} additionally accounts for the current uncertainty
$\bm{\Sigma}$: it focuses sensor resources on state-space directions where
the prior uncertainty is still large.  In the limit $\bm{\Sigma}^{-1} \to
\bm{0}$ (flat prior), the two criteria differ---the BFIM criterion
maximizes $\tr(\bm{F}^\top\bm{F})$, while the A-optimal criterion
minimizes $\tr((\bm{F}^\top\bm{F})^{-1})$.
\end{remark}

\subsection{Stationarity Condition}

At the optimum, the gradient of the full objective vanishes:
\begin{equation}
  \nabla_{\bs}\,\varphi\big|_{\sstar} = \bm{0},
  \qquad\text{i.e.,}\qquad
  \nabla_{\bs}\,\tr(\Snew)\big|_{\sstar} + 2\alpha_r\,\sstar = \bm{0}.
  \label{eq:stationarity}
\end{equation}

The gradient $\nabla_{\bs}\,\tr(\Snew)$ can be computed via the identity
\begin{equation}
  \frac{\partial}{\partial s_k}\,\tr(\bm{M}^{-1})
  = -\tr\!\Bigl(\bm{M}^{-1}\,\frac{\partial \bm{M}}{\partial s_k}\,\bm{M}^{-1}\Bigr),
  \label{eq:tr-inv-deriv}
\end{equation}
where $\bm{M} = \bm{\Sigma}^{-1} + \bm{F}^\top\bm{F}/\sigma^2$ and
\begin{equation}
  \frac{\partial \bm{M}}{\partial s_k}
  = \frac{1}{\sigma^2}\Bigl(
    \frac{\partial \bm{F}^\top}{\partial s_k}\bm{F}
    + \bm{F}^\top \frac{\partial \bm{F}}{\partial s_k}
  \Bigr).
  \label{eq:dM-ds}
\end{equation}

%======================================================================
\section{Outer Optimization and Backpropagation}
\label{sec:outer}
%======================================================================

\subsection{Bilevel Structure}

The outer problem optimizes the sensor design $\bc$ to minimize the
expected estimation error over $N$-step episodes:
\begin{equation}
  \min_{\bc}\;\mathbb{E}_{x_0}\bigl[\|\bmu_N(\bc) - \bx_0\|^2\bigr],
  \label{eq:outer}
\end{equation}
where $\bmu_N$ is the EKF mean after $N$ measurement-update steps, each
involving the sensor selection \eqref{eq:inner-opt}.  The full
computational graph is:
\begin{equation}
  \bc \;\longrightarrow\;
  \underbrace{\sstar_k = \arg\min_{\bs}\,\varphi(\bs;\,\bmu_{k-1},\bm{\Sigma}_{k-1},\bc)}_{\text{inner optimization}}
  \;\longrightarrow\;
  \by_k
  \;\longrightarrow\;
  (\bmu_k, \bm{\Sigma}_k)
  \;\longrightarrow\; \cdots \;\longrightarrow\;
  \|\bmu_N - \bx_0\|^2.
  \label{eq:comp-graph}
\end{equation}

Differentiating through the inner $\arg\min$ requires the Implicit
Function Theorem (IFT).

\subsection{Implicit Function Theorem for $\sstar$}

The stationarity condition \eqref{eq:stationarity} defines $\sstar$
implicitly as a function of $\bm{\theta} \in \{\bc, \bmu, \bm{\Sigma}\}$.
Differentiating with respect to~$\bm{\theta}$:
\begin{equation}
  \bm{H}_\varphi\,\frac{\partial \sstar}{\partial \bm{\theta}}
  + \frac{\partial (\nabla_{\bs}\varphi)}{\partial \bm{\theta}} = \bm{0},
  \label{eq:ift-identity}
\end{equation}
where $\bm{H}_\varphi = \nabla^2_{\bs\bs}\,\varphi\big|_{\sstar}$ is
the Hessian of the full objective at the optimum, which is positive
definite (guaranteed by the regularization $\alpha_r$).

Solving for the sensitivity:
\begin{equation}
  \frac{\partial \sstar}{\partial \bm{\theta}}
  = -\bm{H}_\varphi^{-1}\,
    \frac{\partial (\nabla_{\bs}\varphi)}{\partial \bm{\theta}}
  = \bm{H}_\varphi^{-1}\,\bm{J}_\theta,
  \label{eq:sensitivity}
\end{equation}
where $\bm{J}_\theta = -\partial(\nabla_{\bs}\varphi)/\partial\bm{\theta}$
is the cross-Jacobian.  Since $\nabla_{\bs}\varphi = \nabla_{\bs}\tr(\Snew)
+ 2\alpha_r \bs$ and the regularization term does not depend on $\bm{\theta}$:
\begin{equation}
  \bm{J}_\theta
  = -\frac{\partial}{\partial \bm{\theta}}\,
    \nabla_{\bs}\,\tr(\Snew)\big|_{\sstar}.
\end{equation}

\subsection{Reverse-Mode (Pullback) Computation}

For a scalar loss $\mathcal{L}$ with upstream cotangent
$\bar{\bs} = \partial\mathcal{L}/\partial\sstar \in \R^{d_s}$, the IFT
gives the cotangents for each parameter:
\begin{equation}
  \boxed{%
    \bar{\bm{\theta}} = \bm{J}_\theta^\top\,\bla,
    \qquad\text{where}\quad
    \bm{H}_\varphi\,\bla = \bar{\bs}.
  }
  \label{eq:pullback}
\end{equation}

Concretely, the three cotangents are:
\begin{align}
  \bar{\bc}   &= \bm{J}_c^\top \bla,   &
  \bm{J}_c &= \frac{\partial\,\nabla_{\bs}\tr(\Snew)}{\partial \bc}\in\R^{d_s\times d_c},
  \label{eq:cbar} \\[4pt]
  \bar{\bmu}  &= \bm{J}_\mu^\top \bla,  &
  \bm{J}_\mu &= \frac{\partial\,\nabla_{\bs}\tr(\Snew)}{\partial \bmu}\in\R^{d_s\times d_x},
  \label{eq:mubar} \\[4pt]
  \bar{\bm{\Sigma}}
              &= \operatorname{reshape}(\bm{J}_\Sigma^\top \bla,\; d_x, d_x), &
  \bm{J}_\Sigma &= \frac{\partial\,\nabla_{\bs}\tr(\Snew)}{\partial\ve(\bm{\Sigma})}
  \in\R^{d_s\times d_x^2}.
  \label{eq:Sigmabar}
\end{align}

\begin{remark}[New dependency on $\bm{\Sigma}$]
Unlike the BFIM criterion (where $\sstar$ depends only on $\bmu$ and $\bc$),
the A-optimal policy makes $\sstar$ depend on $\bm{\Sigma}$ as well.
This introduces a new cotangent $\bar{\bm{\Sigma}}$ in the pullback,
which propagates gradients through the evolving covariance across EKF
steps.
\end{remark}

%======================================================================
\section{Efficient Computation of $\bar{\bc}$}
\label{sec:cbar}
%======================================================================

The design-parameter cotangent $\bar{\bc} = \bm{J}_c^\top\bla$ is the
most expensive term because $d_c$ is typically large (e.g., $d_c = 96
n_\omega$ for the photonic sensor design problem, where $n_\omega$ is the
number of frequency points).

\subsection{Na\"ive Approach: Full Jacobian}

Computing $\bm{J}_c$ via forward-mode AD (ForwardDiff) requires
$\lceil d_c/\text{chunk}\rceil$ evaluations of $\nabla_{\bs}\tr(\Snew)$,
each involving nested dual numbers (ForwardDiff-over-ForwardDiff).  For
$d_c = 1920$ and chunk size~8, this is $\sim 240$ passes.

\subsection{Analytical Directional Derivative}

We observe that $\bar{\bc} = \bm{J}_c^\top\bla = \nabla_{\bc}\bigl[
\bla^\top \nabla_{\bs}\tr(\Snew)\bigr]$, which is the gradient of a
\emph{scalar} function of $\bc$.  This scalar is the directional
derivative of $\tr(\Snew)$ in the $\bs$-direction $\bla$:
\begin{equation}
  \bla^\top\nabla_{\bs}\,\tr(\Snew)
  = \frac{\dd}{\dd\epsilon}\,\tr(\Snew(\sstar + \epsilon\bla))
    \Big|_{\epsilon=0}.
  \label{eq:dir-deriv-defn}
\end{equation}

Using \eqref{eq:tr-inv-deriv} with $\bm{M}(\epsilon) = \bm{\Sigma}^{-1}
+ \bm{F}(\sstar+\epsilon\bla)^\top\bm{F}(\sstar+\epsilon\bla)/\sigma^2$:
\begin{equation}
  \frac{\dd}{\dd\epsilon}\,\tr(\bm{M}^{-1})\Big|_{\epsilon=0}
  = -\tr\!\Bigl(\Snew\,\frac{\dd\bm{M}}{\dd\epsilon}\Big|_0\,\Snew\Bigr),
  \label{eq:dir-deriv-trace}
\end{equation}
where
\begin{equation}
  \frac{\dd\bm{M}}{\dd\epsilon}\bigg|_0
  = \frac{1}{\sigma^2}\bigl(
    \dot{\bm{F}}^\top\bm{F} + \bm{F}^\top\dot{\bm{F}}
  \bigr),
  \qquad
  \dot{\bm{F}} \equiv \sum_{k=1}^{d_s} \lambda_k
    \frac{\partial \bm{F}}{\partial s_k}\bigg|_{\sstar}.
  \label{eq:dF-lambda}
\end{equation}

\subsection{Reverse-Mode Gradient}

The key insight is that \eqref{eq:dir-deriv-trace} can be evaluated
\emph{without} any ForwardDiff: $\bm{F}$ and $\dot{\bm{F}}$ are computed
analytically by the lattice scattering model (since the sensor parameters
$\bs$ enter only through the input amplitudes $\bm{a}_{\mathrm{in}}$,
whose $\bs$-derivative is trivial).

The scalar \eqref{eq:dir-deriv-trace} depends on $\bc$ through $\bm{F}$,
$\dot{\bm{F}}$, and $\Snew$ (all of which depend on the S-matrices
packed in $\bc$).  A single reverse-mode AD pass (Zygote) through this
scalar yields $\bar{\bc}$:
\begin{equation}
  \boxed{%
    \bar{\bc} = \nabla_{\bc}\Bigl[
      -\tr\!\bigl(\Snew\,(\dot{\bm{F}}^\top\bm{F}
      + \bm{F}^\top\dot{\bm{F}})\,\Snew/\sigma^2\bigr)
    \Bigr].
  }
  \label{eq:cbar-analytical}
\end{equation}

This reduces the cost from $\mathcal{O}(d_c/\text{chunk})$ nested-dual
passes to \textbf{one} forward + one backward pass through the lattice
model, regardless of~$d_c$.

%======================================================================
\section{Complete Algorithm}
\label{sec:algorithm}
%======================================================================

\begin{algorithm}[h]
\caption{One training iteration (outer gradient step)}
\label{alg:training}
\begin{algorithmic}[1]
\Require Design parameters $\bc$; prior $(\bmu_0, \bm{\Sigma}_0)$;
  true states $\{\bx_0^{(i)}\}$; noise bank $\{\bep_k^{(i)}\}$
\Ensure Loss $\mathcal{L}$ and gradient $\nabla_{\bc}\mathcal{L}$

\Statex \textbf{--- Forward: FDFD simulation ---}
\State $\bc_{\mathrm{flat}} \gets \textsc{SimGeom}(\bm{\varepsilon}_{\mathrm{geom}})$
  \Comment{FDFD $\to$ S-matrices $\to$ flat vector}
\State Store pullback $\overline{\textsc{SimGeom}}$ for later

\Statex \textbf{--- Forward + backward: episode gradients (parallel over episodes) ---}
\For{each episode $i$ \textbf{in parallel}}
  \State $(\bmu, \bm{\Sigma}) \gets (\bmu_0, \bm{\Sigma}_0)$
  \For{$k = 1, \ldots, N$}
    \State $\sstar_k \gets \arg\min_{\bs}\bigl[\tr(\Snew(\bs))
      + \alpha_r\|\bs\|^2\bigr]$
      \Comment{L-BFGS, \Cref{eq:inner-opt}}
    \State $\by_k \gets f(\bx_0^{(i)}, \sstar_k, \bc_{\mathrm{flat}})
      + \bep_k^{(i)}$
    \State $(\bmu, \bm{\Sigma}) \gets \textsc{EKF-Update}(\bmu, \bm{\Sigma},
      \by_k, \sstar_k, \bc_{\mathrm{flat}})$
  \EndFor
  \State $\mathcal{L}^{(i)} \gets \|\bmu - \bx_0^{(i)}\|^2$
  \State Backpropagate through EKF loop using IFT (\Cref{eq:pullback})
    at each $\arg\min$
  \State \Return $(\mathcal{L}^{(i)},\; \nabla_{\bc_{\mathrm{flat}}}
    \mathcal{L}^{(i)})$
\EndFor

\Statex \textbf{--- Backward: FDFD adjoint ---}
\State $\bar{\bc}_{\mathrm{flat}} \gets
  \frac{1}{n_{\mathrm{ep}}}\sum_i \nabla_{\bc_{\mathrm{flat}}}
  \mathcal{L}^{(i)}$
\State $\nabla_{\bm{\varepsilon}_{\mathrm{geom}}}\mathcal{L}
  \gets \overline{\textsc{SimGeom}}(\bar{\bc}_{\mathrm{flat}})$
  \Comment{Adjoint of FDFD solver}

\Statex \textbf{--- Adam update ---}
\State $\bm{\varepsilon}_{\mathrm{geom}} \gets
  \textsc{Adam-Step}(\bm{\varepsilon}_{\mathrm{geom}},\;
  \nabla_{\bm{\varepsilon}_{\mathrm{geom}}}\mathcal{L})$
\end{algorithmic}
\end{algorithm}

%======================================================================
\section{Summary of Derivatives}
\label{sec:derivatives}
%======================================================================

\Cref{tab:derivatives} summarizes all derivative quantities needed by the
algorithm, their dimensionality, and the method used to compute them.

\begin{table}[h]
\centering
\caption{Derivative quantities in the bilevel optimization.}
\label{tab:derivatives}
\begin{tabular}{@{}llll@{}}
\toprule
\textbf{Quantity} & \textbf{Size} & \textbf{Method} & \textbf{Used in} \\
\midrule
$\nabla_{\bs}\tr(\Snew)$ & $d_s$ & ForwardDiff & Inner L-BFGS \\
$\nabla^2_{\bs\bs}\tr(\Snew)$ & $d_s \times d_s$ & ForwardDiff (nested) & IFT Hessian $\bm{H}_\varphi$ \\
$\bla = \bm{H}_\varphi^{-1}\bar{\bs}$ & $d_s$ & LU solve & IFT pullback \\
$\bm{J}_c^\top\bla = \bar{\bc}$ & $d_c$ & Zygote (analytical) & Outer gradient \\
$\bm{J}_\mu^\top\bla = \bar{\bmu}$ & $d_x$ & ForwardDiff & EKF chain rule \\
$\bm{J}_\Sigma^\top\bla = \bar{\bm{\Sigma}}$ & $d_x \times d_x$ & ForwardDiff & EKF chain rule \\
$\bm{F}$, $\dot{\bm{F}}$ & $d_y \times d_x$ & Analytical (lattice) & Eq.~\eqref{eq:cbar-analytical} \\
$\Snew$ & $d_x \times d_x$ & Woodbury \eqref{eq:woodbury} & Eq.~\eqref{eq:cbar-analytical} \\
\bottomrule
\end{tabular}
\end{table}

\end{document}
